\chapter{Reasoning Frameworks}
\label{ch:reasoning}

\epigraph{``Reasoning is the bridge between perception and action. The quality of the bridge determines the quality of the agent.''}{}

\learninggoal{
	\item Distinguish between single-pass reasoning, iterative refinement, and search-based reasoning.
	\item Implement ReAct, Chain-of-Thought, and Tree-of-Thought patterns.
	\item Understand the role of reflection as verbal reinforcement learning.
	\item Compare reasoning frameworks along axes of compute cost, reliability, and generality.
}

\section{Why Reasoning Matters}

A bare \llm\ generates a single continuation from a prompt. An agent needs structured reasoning because:

\begin{itemize}
	\item \textbf{Sequential dependencies:} earlier actions constrain later ones. The agent must reason about the sequence, not just the next token.
	\item \textbf{Exploration:} the best action is not always the highest-probability one. The agent must consider alternatives.
	\item \textbf{Error recovery:} actions fail. The agent must detect failure and adjust.
	\item \textbf{Multi-step tasks:} complex goals decompose into sub-goals that require intermediate reasoning.
\end{itemize}

\begin{definition}[Reasoning Framework]
	A \textbf{reasoning framework} is a structured prompting strategy that elicits multi-step deliberation from an \llm, typically by interleaving reasoning traces (\emph{thoughts}) with actions or intermediate conclusions.
\end{definition}

\section{Chain-of-Thought (CoT)}

The simplest structured reasoning method~\cite{wei2022cot}. The \llm\ is prompted to produce intermediate reasoning steps before the final answer.

\begin{lstlisting}[caption=Chain-of-Thought prompt structure]
  System: Solve the problem step by step.

  User: If a store has 15 apples and sells 3 each day,
        how many days until they have 6 left?

  Assistant: Let me think step by step.
  1. Starting apples: 15
  2. Target apples remaining: 6
  3. Apples that need to be sold: 15 - 6 = 9
  4. They sell 3 per day, so days needed: 9 / 3 = 3
  Answer: 3 days.
\end{lstlisting}

CoT is zero-shot (``let's think step by step'') or few-shot (provide examples). It improves performance on arithmetic, logical, and commonsense reasoning tasks.

\begin{theorem}[CoT Improves Expressiveness]
	There exist problems that a transformer of depth $L$ cannot solve without CoT but can solve with $O(n)$ CoT steps, because CoT increases the effective circuit depth available for computation (see the companion volume on LLM theory).
\end{theorem}

\subsection{Self-Consistency}

A natural extension: sample multiple CoT paths and take the majority answer~\cite{wang2023selfconsistency}. This smooths out the variance of individual CoT samples.

\begin{equation}
	\hat{y} = \text{majority-vote}\left(\bigcup_{i=1}^k \text{CoT}(x; \theta_i)\right)
\end{equation}

With $k=5$ to $k=20$ samples, accuracy improves by 5--15\% on reasoning benchmarks at the cost of $k \times$ compute.

\section{ReAct: Reasoning + Acting}

ReAct~\cite{yao2023react} interleaves reasoning traces (\emph{thoughts}) with actions. The reasoning trace serves to:
\begin{enumerate}
	\item Track the current state of the task.
	\item Decompose the next step into a reasoned decision.
	\item Justify why a particular action was chosen.
\end{enumerate}

\begin{definition}[ReAct Trajectory]
	A ReAct trajectory is a sequence of interleaved \textbf{thoughts}, \textbf{actions}, and \textbf{observations}:
	\begin{equation}
		\tau = (t_1, a_1, o_1, t_2, a_2, o_2, \dots, t_k, a_k, o_k, \text{answer})
	\end{equation}
	where $t_i$ is a reasoning step, $a_i$ is a tool call, and $o_i$ is the tool's output.
\end{definition}

\begin{example}[ReAct in Practice]
	\begin{lstlisting}
  Thought: I need to find the current population of Tokyo.
  Action: search("Tokyo population 2025")
  Observation: "Tokyo has 14.1 million residents (2025 estimate)"
  Thought: The user also asked about the area.
  Action: search("Tokyo area square kilometers")
  Observation: "Tokyo covers 2,194 km²"
  Thought: Now I can compute density.
  Action: calculator(14.1e6 / 2194)
  Observation: "6426"
  Answer: Tokyo has approximately 6,426 people per km².
  \end{lstlisting}
\end{example}

\section{Tree of Thoughts (ToT)}

ToT~\cite{yao2023tree} generalises CoT from a single chain to a tree of reasoning paths, using search to explore alternatives.

\begin{definition}[Tree of Thoughts]
	ToT maintains a tree of \textbf{thoughts}---intermediate reasoning states---and uses search algorithms (BFS, DFS, or MCTS) to explore the tree. At each node, the \llm\ either:
	\begin{itemize}
		\item \textbf{Generates} successor thoughts (branching), or
		\item \textbf{Evaluates} the promise of a thought (pruning).
	\end{itemize}
\end{definition}

\begin{algorithm}[ToT with BFS]
	\begin{enumerate}
		\item Initialise the tree with the input $x$ as the root.
		\item For each level $d = 1, \dots, D$:
		      \begin{enumerate}
			      \item For each node at depth $d-1$, prompt the \llm\ to generate $b$ candidate thoughts.
			      \item For each candidate, prompt the \llm\ to evaluate its promise (a score or classification).
			      \item Retain the top $k$ most promising nodes.
		      \end{enumerate}
		\item Select the path from root to the highest-scoring leaf as the reasoning chain.
	\end{enumerate}
\end{algorithm}

\begin{example}[ToT for Game of 24]
	\begin{lstlisting}
  Input: 4, 7, 8, 8
  Level 1 thoughts:
    - "4 + 7 = 11" → remaining: 11, 8, 8  (score: 0.6)
    - "8 - 4 = 4"  → remaining: 4, 7, 8   (score: 0.4)
    - "8 * 7 = 56" → remaining: 56, 4, 8  (score: 0.2)
  Prune: keep top 2.
  Level 2 from "4 + 7 = 11":
    - "11 + 8 = 19" → remaining: 19, 8    (score: 0.7)
    - "11 - 8 = 3"  → remaining: 3, 8     (score: 0.5)
  ... continue until (8 - 7) * 4 + 8 = 12  -> wait that's wrong
  \end{lstlisting}
\end{example}

Cost: ToT using BFS with branching factor $b=5$, depth $D=3$, and top-$k=2$ requires up to $1 + 5 + 10 = 16$ \llm\ calls per problem, compared to 1 for CoT.

\section{Reflection and Reflexion}

Reflection adds an evaluation step after actions, generating verbal feedback that improves future actions.

\subsection{Reflexion}

Reflexion~\cite{shinn2023reflexion} maintains a \textbf{reflection memory} that stores natural-language summaries of past failures:

\begin{lstlisting}[caption=Reflexion architecture]
  # After each episode:
  trajectory = [(thought_1, action_1, observation_1), ...]
  reflection = llm(f"Review this trajectory. What went wrong?
                    What should I do differently next time?
                    Trajectory: {trajectory}")
  memory.append(reflection)

  # On the next episode:
  system_prompt = base_prompt + "\nLessons from past attempts:\n"
                  + "\n".join(memory)
\end{lstlisting}

The reflection serves as a cheap substitute for gradient updates: the agent ``learns'' from mistakes without updating model weights.

\begin{theorem}[Reflection as Verbal RL]
	Reflexion approximates policy improvement by transforming environment feedback into natural-language instructions that condition future behaviour. This is ``verbal reinforcement learning''---the policy is updated by modifying the prompt rather than the weights.
\end{theorem}

\subsection{Comparison with RL Fine-Tuning}

\begin{longtable}{p{3.5cm}p{4cm}p{4cm}p{3cm}}
	\toprule
	\textbf{Property}  & \textbf{Reflexion}            & \textbf{RL Fine-Tuning} & \textbf{Best for} \\
	\midrule
	Weight update      & None                          & Full model              & --                \\
	Sample efficiency  & 1 episode                     & $10^4$--$10^6$ episodes & Reflexion         \\
	Generalisation     & Limited to similar tasks      & Broad                   & RL                \\
	Compute per update & 1 \llm\ call                  & Full training run       & Reflexion         \\
	Stability          & Unstable (prompt brittleness) & Stable with tuning      & RL                \\
	\bottomrule
\end{longtable}

\section{Choosing a Reasoning Framework}

The right framework depends on the task's structure and the cost of mistakes:

\begin{longtable}{p{3cm}p{4cm}p{4cm}p{4cm}}
	\toprule
	\textbf{Task type}       & \textbf{Recommended} & \textbf{Why}                   & \textbf{Compute cost} \\
	\midrule
	Simple QA                & CoT                  & Enough reasoning, minimal cost & 1 pass                \\
	Multi-step search        & ReAct                & Interleaved reasoning + tools  & 3--10 passes          \\
	Creative problem solving & ToT                  & Exploration of alternatives    & 10--100 passes        \\
	Error-prone tasks        & Reflexion            & Iterative improvement          & 2--5 episodes         \\
	Open-ended exploration   & Hierarchical         & Decomposition + focused search & 10--50 passes         \\
	\bottomrule
\end{longtable}

\section{Exercises}

\begin{exercise}
	Implement ReAct for a simple tool set (calculator, date lookup, Wikipedia search). Show the full trajectory for a question that requires 2--3 tool calls.
\end{exercise}

\begin{exercise}
	Compare CoT and ToT on a logic puzzle (e.g., Einstein's zebra puzzle). How many \llm\ calls does each method require? Which produces the correct answer?
\end{exercise}

\begin{exercise}
	Design a reflection prompt for a coding agent. What information should it include about the failed attempt? How would you structure the memory to avoid catastrophic forgetting of useful lessons?
\end{exercise}

\begin{exercise}
	Prove that the expected number of \llm\ calls in ToT with BFS, branching factor $b$, depth $D$, and top-$k$ pruning is $1 + D \cdot b \cdot k$. Generalise to DFS.
\end{exercise}
